% Gesamttest «Dreiecke» — Quelle des PDFs (python3 scripts/build-lp-pdf.py leitprogramme/dreiecke/).
% Musterlösung und Raster: bewertungspaket.tex. Alle Werte mit python3 nachgerechnet (scripts/lp/dreiecke/zahlen.py, 08.10.2026).
% Nach der Prüfung vom 08.10.2026 überarbeitet: G2 ohne (d), G4 neue Zahlen und (c) Abstand, G7 Schwerpunkt statt Jonas.
% Ein Teil: RLP 5.2 trägt keinen Vermerk «auch ohne Hilfsmittel», der Taschenrechner ist überall erlaubt.
\documentclass[11pt]{article}
\usepackage{../lp-druck}
\renewcommand{\lpname}{Dreiecke}
\renewcommand{\lpteilgebiet}{5.2a}
\renewcommand{\lpbereich}{Grundlagenfach}
\renewcommand{\lpfuss}{Gesamttest · 25 Punkte}
\hyphenation{Mittel-senkrechte Mittel-senkrechten Winkel-halbierende Winkel-halbierenden Seiten-halbierende Seiten-halbierenden}
\begin{document}
\lpkopf{Gesamttest}{%
  \vspace{4pt}\noindent\begin{tabularx}{\linewidth}{@{}X X@{}}
  Code (kein Name): \hrulefill & Punkte: \hrulefill\ / 25
  \end{tabularx}}

\begin{lpinfo}{Anleitung}
\textbf{Zeit:} rund 30 Minuten. \textbf{Hilfsmittel:} Taschenrechner erlaubt (RLP 5.2 trägt keinen Vermerk «auch ohne Hilfsmittel»).
Mach zu jeder Aufgabe eine \textbf{Skizze} und schreib den \textbf{Rechenweg} auf — bewertet wird der Weg. Winkel in Grad,
Längen auf zwei Dezimalen runden, Zwischenresultate ungerundet weiterverwenden. Figuren sind Skizzen, nicht massstäblich.
Danach: Lösung scannen oder fotografieren (ohne Namen und Standort) und mit dem \emph{Bewertungspaket} bewerten lassen.
\end{lpinfo}

\lpteil{Kapitel 1 bis 4}{Taschenrechner erlaubt · 25 Punkte}

\begin{lpaufgabe}{G1}{Winkel aus dem Aussenwinkel}{4 P}
In einem Dreieck ist $\beta$ doppelt so gross wie $\alpha$, und der Aussenwinkel bei $C$ misst $126^\circ$.
\begin{enumerate}[label=(\alph*),leftmargin=*,itemsep=2pt]
\item Berechne $\alpha$, $\beta$ und $\gamma$.
\item Ist das Dreieck spitz-, recht- oder stumpfwinklig? Ist es gleichschenklig? Begründe.
\end{enumerate}
\lpplatz{6}
\end{lpaufgabe}

\begin{lpaufgabe}{G2}{Brunnen und Bank im Park}{3 P}
Ein Park hat die Form eines Dreiecks $ABC$; seine Ränder sind gerade Wege.
\begin{enumerate}[label=(\alph*),leftmargin=*,itemsep=2pt]
\item Ein Brunnen soll von allen drei Ecken gleich weit entfernt sein. Welche Linien zeichnest du, um seinen Ort zu finden? Wie heisst der Punkt?
\item Der Park hat bei $C$ einen Winkel von $105^\circ$. Liegt der Brunnen im Park? Begründe.
\item Eine Bank soll von allen drei Wegen gleich weit entfernt sein. Welche Linien, welcher Punkt?
\end{enumerate}
\lpplatz{6}
\end{lpaufgabe}

\begin{lpaufgabe}{G3}{Höhe und Winkelhalbierende}{4 P}
Im Dreieck $ABC$ ist $\alpha = 74^\circ$ und $\beta = 46^\circ$. Die Höhe $h_c$ trifft die Seite $c$ in $F$, die Winkelhalbierende
$w_\gamma$ trifft $c$ in $D$. Wie gross ist der Winkel zwischen $h_c$ und $w_\gamma$ bei $C$? Skizziere und begründe jeden Schritt.
\lpplatz{8}
\end{lpaufgabe}

\begin{lpaufgabe}{G4}{Eine fremde Lösung}{4 P}
\begin{minipage}[t]{0.50\linewidth}\vspace{0pt}
Im Dreieck $ABC$ ist $c = \overline{AB} = 7\,\text{cm}$, $a = \overline{BC} = 5.8\,\text{cm}$ und die Höhe $h_c = 4.2\,\text{cm}$.
Ihr Fusspunkt $F$ liegt auf der Verlängerung von $c$, $4\,\text{cm}$ hinter $B$. Mia rechnet:
\[ A = \tfrac{1}{2} \cdot 7 \cdot 5.8 = 20.3\,\text{cm}^2 \]
\begin{enumerate}[label=(\alph*),leftmargin=*,itemsep=2pt]
\item Was hat Mia falsch gemacht? Berechne die Fläche richtig.
\item Berechne die Seite $b = \overline{AC}$.
\item Wie weit ist die Ecke $B$ von der Geraden $AC$ entfernt?
\end{enumerate}
\end{minipage}\hfill
\begin{minipage}[t]{0.46\linewidth}\vspace{0pt}\centering
\begin{tikzpicture}[scale=0.40]
  \coordinate (A) at (0,0); \coordinate (B) at (7,0); \coordinate (C) at (11,4.2); \coordinate (F) at (11,0);
  \draw[lpblau,thick,fill=lpblau!8] (A)--(B)--(C)--cycle;
  \draw[lpgrau,dashed] (B)--(F);
  \draw[lporange,thick,dashed] (C)--(F); \draw[lporange] (11,0) rectangle (10.6,0.4);
  \node[below left] at (A) {$A$}; \node[below] at (B) {$B$}; \node[above] at (C) {$C$}; \node[below right] at (F) {$F$};
  \node[below] at (3.5,0) {\small $7\,\text{cm}$}; \path (B) -- (C) node[midway, sloped, below] {\small $5.8\,\text{cm}$};
  \node[right] at (11,2.1) {\small $h_c = 4.2\,\text{cm}$}; \node[below] at (9,0) {\small $4\,\text{cm}$};
\end{tikzpicture}
\end{minipage}
\lpplatz{6}
\end{lpaufgabe}

\begin{lpaufgabe}{G5}{Rechtwinkliges Dreieck in schräger Lage}{3 P}
\begin{minipage}[t]{0.60\linewidth}\vspace{0pt}
Im Dreieck $KLM$ liegt der rechte Winkel bei $L$; $\overline{KL} = 4\,\text{cm}$, $\overline{KM} = 8.5\,\text{cm}$.
\begin{enumerate}[label=(\alph*),leftmargin=*,itemsep=2pt]
\item Welche Seite ist die Hypotenuse?
\item Berechne $\overline{LM}$.
\item Berechne die Fläche des Dreiecks.
\end{enumerate}
\end{minipage}\hfill
\begin{minipage}[t]{0.36\linewidth}\vspace{0pt}\centering
\begin{tikzpicture}[scale=0.42]
  \coordinate (K) at (0,0); \coordinate (L) at (3.625,1.690); \coordinate (M) at (0.455,8.487);
  \draw[lpblau,thick,fill=lpblau!8] (K)--(L)--(M)--cycle;
  \draw (3.172,1.479) -- (2.961,1.932) -- (3.414,2.143);
  \node[below left] at (K) {$K$}; \node[right] at (L) {$L$}; \node[above] at (M) {$M$};
  \node[below right] at (1.8,0.85) {\small $4\,\text{cm}$}; \node[left] at (0.2,4.3) {\small $8.5\,\text{cm}$};
\end{tikzpicture}
\end{minipage}
\lpplatz{5}
\end{lpaufgabe}

\begin{lpaufgabe}{G6}{Giebelwand}{4 P}
Die Giebelwand eines Hauses ist ein gleichschenkliges Dreieck: unten $9.6\,\text{m}$ breit, der First liegt $3.6\,\text{m}$ über der Grundlinie.
\begin{enumerate}[label=(\alph*),leftmargin=*,itemsep=2pt]
\item Wie lang ist eine Dachschräge (vom Rand der Grundlinie bis zum First)?
\item Wie gross ist die Fläche der Giebelwand?
\item Wie lang ist der Umfang des Giebeldreiecks?
\end{enumerate}
\lpplatz{6}
\end{lpaufgabe}

\begin{lpaufgabe}{G7}{Ein Blech balancieren}{3 P}
\begin{minipage}[t]{0.56\linewidth}\vspace{0pt}
Ein dreieckiges Blech $ABC$ hat bei $C$ einen rechten Winkel; $\overline{CA} = 7\,\text{cm}$, $\overline{CB} = 12\,\text{cm}$.
Im Schwerpunkt $S$ lässt es sich auf einer Nadelspitze balancieren. $M_a$ ist die Mitte von $\overline{BC}$.
\begin{enumerate}[label=(\alph*),leftmargin=*,itemsep=2pt]
\item Wie lang ist die Seitenhalbierende $s_a = \overline{AM_a}$?
\item Wie weit ist $S$ von $A$ entfernt?
\end{enumerate}
\end{minipage}\hfill
\begin{minipage}[t]{0.40\linewidth}\vspace{0pt}\centering
\begin{tikzpicture}[scale=0.32]
  \coordinate (C) at (0,0); \coordinate (B) at (12,0); \coordinate (A) at (0,7); \coordinate (M) at (6,0);
  \draw[lpblau,thick,fill=lpblau!8] (A)--(B)--(C)--cycle;
  \draw (0,0.6) -- (0.6,0.6) -- (0.6,0);
  \draw[lporange,thick,dashed] (A)--(M); \fill (M) circle (2.5pt);
  \node[left] at (A) {$A$}; \node[below right] at (B) {$B$}; \node[below left] at (C) {$C$}; \node[below] at (M) {$M_a$};
\end{tikzpicture}
\end{minipage}
\lpplatz{5}
\end{lpaufgabe}
\end{document}
